% ============================================================================
%  part2a_inference.tex  --  Part II, sections 4-5.
%  Sparse variational GP inference + the eigenspace representation and the two
%  current implementations (vargp_direct / default_gpy).
%  Self-contained \input-ready fragment; uses ONLY the shared preamble macros.
%  Ground truth = current code; where a doc disagreed, the code wins (memo 4).
% ============================================================================

% Scoped \sloppy so long code identifiers (file.py:function) wrap instead of
% bleeding into the margin; the \begingroup/\endgroup confines it to this
% fragment so no global typesetting state leaks into neighbouring sections.
\begingroup\sloppy

\section{Sparse variational Gaussian-process inference}\label{sec:varinf}

Part~I fixed the generative model: a zero-mean latent Gaussian process
$\lat \sim \GProc(0,\Kf)$ over image space, an exponential link to the firing
rate $\rate(x)=\exp\!\big(\gain\,\lat(x)+\latz\big)$, and Poisson spike counts
$\spk_i \sim \Poi\!\big(\rate(x_i)\big)$, where $\Kf$ is the first-order
arc-cosine kernel of \eqref{eq:acoskernel} built on the structured spatial
prior covariance $\Cmat$ of \eqref{eq:Cprior}. Here we turn that model into a
trainable posterior. Two current implementations realise \emph{the same}
inference; both are introduced abstractly in this section and made concrete in
\S\ref{sec:eigenspace}.

\subsection{Poisson breaks conjugacy}\label{sec:conjugacy}

Given a batch of images $X=(x_1,\dots,x_{\Ntr})$ and spike counts
$\spk=(\spk_1,\dots,\spk_{\Ntr})$, the exact posterior over the latent and the
marginal likelihood,
\begin{align}
  p(\lat\mid\spk,X,\theta)&\;\propto\;
    \Big[\textstyle\prod_i \Poi\!\big(\spk_i\mid e^{\gain\lat_i+\latz}\big)\Big]\,
    \Gauss(\lat\mid 0,\Kf_\theta), \notag\\
  p(\spk\mid X,\theta)&=\!\int\! \prod_i \Poi\!\big(\spk_i\mid e^{\gain\lat_i+\latz}\big)\,
    \Gauss(\lat\mid 0,\Kf_\theta)\,d\lat ,
\end{align}
have no closed form: the exponential-Poisson likelihood is not conjugate to the
Gaussian prior, so the integral over $\lat$ is intractable. We therefore
approximate the posterior by a tractable Gaussian and maximise a variational
lower bound (\S\ref{sec:elbo}).

\begin{remark}[The conjugate contrast]
For a \emph{Gaussian} observation $\spk=\lat+\varepsilon$,
$\varepsilon\sim\Gauss(0,\pvar_r I)$, everything would be closed-form: the
marginal is $\spk\mid X\sim\Gauss(0,\Kf+\pvar_r I)$, the predictive at a test
image $x_\ast$ is Gaussian with mean $\kvec_\ast\transpose(\Kf+\pvar_r I)^{-1}\spk$
and variance $\Kf(x_\ast,x_\ast)-\kvec_\ast\transpose(\Kf+\pvar_r I)^{-1}\kvec_\ast$,
and the log-evidence is
$-\tfrac12\spk\transpose(\Kf+\pvar_r I)^{-1}\spk-\tfrac12\log\lvert\Kf+\pvar_r I\rvert-\tfrac{\Ntr}{2}\log 2\pi$.
None of these survive the Poisson link; the variational route below recovers a
usable posterior in its place.
\end{remark}

\subsection{Inducing points and the variational posterior}\label{sec:inducing}

\begin{definition}[Sparse variational posterior]
Introduce $\Mind$ \emph{inducing points} $\Ipt=\{\ipt_1,\dots,\ipt_{\Mind}\}$
with latent values $\latt=\lat(\Ipt)$ (the code's inducing values $u$), and place a
Gaussian variational posterior on them,
\begin{equation}
  q(\latt)=\Gauss(\latt\mid\vm,\vV),\qquad
  \vm\in\Reals^{\Mind},\quad \vV\in\mathbb{S}_{++}^{\Mind}.
  \label{eq:qdef}
\end{equation}
The posterior at \emph{any} image $x$ is recovered by marginalising the
conditional prior $p(\lat(x)\mid\latt)$ against $q(\latt)$.
\end{definition}

Exact inference would need the $\Ntr\!\times\!\Ntr$ kernel over all training
images, an $O(\Ntr^3)$ cost; the inducing construction reduces it to
$O(\Ntr\Mind^2)$.

\begin{remark}[Closed-loop invariant $\Mind=\Ntr$]\label{rem:MeqN}
In the active-learning loop the inducing set is taken equal to the training set,
$X=\Ipt$ (code \texttt{X\_train == X\_tilde}), so $\Mind=\Ntr$ and $\Kf=\Ktil$.
The ``sparse'' posterior is then \emph{exact}. The SVGP formalism is retained not
for sparsity but for two other reasons: (i)~it supplies the closed-form
variational machinery of \S\ref{sec:predmoments}--\S\ref{sec:fbar}, and (ii)~its
inducing structure admits the $O(\Mind)$ rank-one online growth
(\S\ref{sec:rank1}) that appends one image per active step.
\end{remark}

\subsection{Projected posterior moments}\label{sec:predmoments}

Marginalising $p(\lat(x)\mid\latt)\,q(\latt)$ yields a Gaussian
$q(\lat(x))=\Gauss\!\big(\pmean(x),\pvar(x)\big)$ whose moments are the
inducing-point moments projected through the kernel:
\begin{align}
  \pmean(x)  &= \kvec(x)\transpose\,\Ktil^{-1}\,\vm , \notag\\
  \pvar(x)   &= \Kf(x,x) + \kvec(x)\transpose\,\Ktil^{-1}\big(\vV-\Ktil\big)\Ktil^{-1}\,\kvec(x),
  \label{eq:predmoments}
\end{align}
with the cross-covariance vector $\kvec(x)=\Kf(\Ipt,x)\in\Reals^{\Mind}$, the
inducing gram $\Ktil=\Kf(\Ipt,\Ipt)\in\Reals^{\Mind\times\Mind}$, and the prior
self-variance $\Kf(x,x)$ given by the self-kernel identity \eqref{eq:selfkernel}.
Introducing the \emph{projection vector} $\pv(x)=\Ktil^{-1}\kvec(x)$ and the
Schur complement $s(x)=\Kf(x,x)-\kvec(x)\transpose\Ktil^{-1}\kvec(x)$, the
variance splits into a data-independent residual and an epistemic term,
\begin{equation}
  \pvar(x)=\underbrace{\big[\Kf(x,x)-\kvec(x)\transpose\Ktil^{-1}\kvec(x)\big]}_{\text{residual (Nystr\"om/Schur) }s(x)}
    \;+\;\underbrace{\kvec(x)\transpose\Ktil^{-1}\vV\Ktil^{-1}\kvec(x)}_{\text{epistemic}}
    \;=\; s(x)+\pv(x)\transpose\vV\,\pv(x),
  \label{eq:varsplit}
\end{equation}
and the mean reads $\pmean(x)=\pv(x)\transpose\vm$. The two forms of $\pvar$ are
algebraically identical since
$\kvec\transpose\Ktil^{-1}(\vV-\Ktil)\Ktil^{-1}\kvec=\kvec\transpose\Ktil^{-1}\vV\Ktil^{-1}\kvec-\kvec\transpose\Ktil^{-1}\kvec$.

\begin{remark}[Prior sanity check]
At the prior $\vV=\Ktil$ the correction in \eqref{eq:predmoments} vanishes and
$\pvar(x)=\Kf(x,x)$ (the prior variance); as data drive $\vV$ below $\Ktil$,
$\pvar(x)<\Kf(x,x)$.
\end{remark}

Both code paths compute exactly \eqref{eq:predmoments}:
\texttt{vargp\_direct} in the eigenbasis of $\Ktil$
(\texttt{eigenspace\_model.py:\allowbreak \_lambda\_moments\_eigenspace}, \S\ref{sec:eig-moments}),
\texttt{default\_gpy} through GPyTorch's \texttt{VariationalStrategy} internals.

\subsection{The ELBO}\label{sec:elbo}

All parameters are learned by maximising the evidence lower bound on
$\log p(\spk)$,
\begin{equation}
  \ELBO \;=\; \sum_i \E_{q(\lat_i)}\!\big[\log \Poi(\spk_i\mid\lat_i)\big]
    \;-\; \KLdiv\!\big(q(\latt)\,\big\|\,p(\latt)\big),
  \label{eq:elbo}
\end{equation}
which is tight when $q$ equals the true posterior. Because the likelihood
factorises over images and $q(\lat_i)=\Gauss(\pmean_i,\pvar_i)$ is Gaussian, each
expected log-likelihood term is closed-form:
\begin{equation}
  \E_{q(\lat_i)}\!\big[\log \Poi(\spk_i\mid\lat_i)\big]
    \;=\; \spk_i\big(\gain\pmean_i+\latz\big)
      \;-\; \exp\!\big(\gain\pmean_i+\tfrac12\gain^2\pvar_i+\latz\big)
      \;+\; \text{const},
  \label{eq:ell}
\end{equation}
with $\text{const}=-\log(\spk_i!)$ independent of all parameters (dropped in
code). \texttt{likelihoods.py:\allowbreak expected\_log\_prob} returns, summed over images,
\begin{lstlisting}
log_prob = target*(A*mu + lambda0) \
           - exp(A*mu + 0.5*A**2*var + lambda0)
\end{lstlisting}
matching \eqref{eq:ell} term by term. The KL term is the standard
Gaussian--Gaussian divergence between $q(\latt)=\Gauss(\vm,\vV)$ and the prior
$p(\latt)=\Gauss(0,\Ktil)$,
\begin{equation}
  \KLdiv(q\,\|\,p) \;=\; \tfrac12\Big[\,
    \log\frac{\lvert\Ktil\rvert}{\lvert\vV\rvert}
    \;+\; \Tr\!\big(\Ktil^{-1}\vV\big)
    \;+\; \vm\transpose\Ktil^{-1}\vm
    \;-\; \Mind \,\Big].
  \label{eq:kl}
\end{equation}

\begin{remark}[The $-\Mind$ constant]
The $-\Mind$ in \eqref{eq:kl} is parameter-independent (it is the dimension of
the inducing Gaussian) and carries zero gradient, so it affects only the reported
loss magnitude, never the optimum or the predictions. The eigenspace evaluation
of \S\ref{sec:eig-kl} carries the analogous $-\nb$ for the same reason.
\end{remark}

The full objective is maximised by EM-style coordinate ascent
(E-step on $(\vm,\vV)$, F-step on $(\gain,\latz)$, M-step on the kernel
hyperparameters); the update equations are derived in the training-algorithm
section. Equation~\eqref{eq:elbo} is the single objective all three blocks share.

\subsection{Expected firing rate via the Gaussian MGF}\label{sec:fbar}

The exponential term in \eqref{eq:ell} is the \emph{expected firing rate} under
$q$,
\begin{equation}
  \fbar_i \;=\; \E_{q(\lat_i)}\!\big[\exp(\gain\lat_i+\latz)\big]
    \;=\; \exp\!\big(\gain\pmean_i+\tfrac12\gain^2\pvar_i+\latz\big).
  \label{eq:fbar}
\end{equation}
This is the Gaussian moment-generating function
$\E[e^{tZ}]=\exp(t\mu+\tfrac12 t^2\sigma^2)$ applied to $Z=\gain\lat_i+\latz$ at
$t=1$ (code \texttt{likelihoods.py:\allowbreak expected\_firing\_rate}, identical to the
loop's \texttt{compute\_f\_mean}).

\begin{remark}[Jensen inflation]
The term $\tfrac12\gain^2\pvar_i$ is strictly positive whenever $\pvar_i>0$, so
posterior \emph{uncertainty raises} the predicted rate above the point estimate
$\exp(\gain\pmean_i+\latz)$ --- a direct consequence of Jensen's inequality for
the convex map $\exp$. Confidence and rate are thus coupled: a poorly-constrained
image is predicted to fire \emph{more}, not less.
\end{remark}

\subsection{Prediction: expected spike count}\label{sec:prediction}

The moments \eqref{eq:predmoments} do not depend on the observed $\spk$, so the
same formulas serve as the predictive distribution at any novel image $x_\ast$.
The expected spike count is the expected firing rate \eqref{eq:fbar} evaluated
there,
\begin{equation}
  \langle \spk_\ast\rangle=\langle\rate(x_\ast)\rangle
    =\exp\!\big(\gain\pmean(x_\ast)+\tfrac12\gain^2\pvar(x_\ast)+\latz\big),
\end{equation}
with $\pmean(x_\ast),\pvar(x_\ast)$ from \eqref{eq:predmoments} using freshly
evaluated $\kvec(x_\ast)$ and $\Kf(x_\ast,x_\ast)$. The firing rate is
log-normal, so its variance is
\begin{equation}
  \operatorname{Var}\!\big(\rate(x)\big)
    =\langle\rate(x)\rangle^2\big(e^{\gain^2\pvar(x)}-1\big).
  \label{eq:fvar}
\end{equation}

\subsection{The posterior cross-covariance pitfall}\label{sec:crosscov}

The self-variance of \eqref{eq:predmoments} extends formally to the
\emph{cross-covariance} between two distinct images $x$ (a candidate sample
location) and $x_\ast$ (a query location):
\begin{equation}
  \Sigma_{x,x_\ast}=\operatorname{Cov}\!\big[\lat(x),\lat(x_\ast)\mid \mathcal D\big]
    =\Kf(x,x_\ast)+\pv(x)\transpose\big(\vV-\Ktil\big)\pv(x_\ast),
  \qquad \pv(\cdot)=\Ktil^{-1}\kvec(\cdot),
  \label{eq:crosscov}
\end{equation}
which reduces to $\pvar(x)$ at $x=x_\ast$.

\begin{remark}[Failure outside the inducing hull]
For a query $x_\ast$ \emph{outside} the inducing region, \eqref{eq:crosscov}
produces cross-covariances that violate the Cauchy--Schwarz bound
$\lvert\Sigma_{x,x_\ast}\rvert\le\sqrt{\Sigma_{x,x}\,\Sigma_{x_\ast,x_\ast}}$.
Measured example (RBF, $\ell=0.2$, $20$ inducing points in $[-1,1]$, sample
$x=0$): at $x_\ast=\pm1.5$ the manual formula gives $\Sigma\approx-4$ while
GPyTorch's covariance gives $\approx 0$; with self-variances $\sim0.04$ the valid
maximum is $\sim0.04$, so $-4$ is two orders of magnitude out of range. Root
cause: in the residual $\Kf(x,x_\ast)-\pv(x)\transpose\Ktil\,\pv(x_\ast)$ the
direct kernel $\Kf(x,x_\ast)\!\to\!0$ (finite lengthscale) outside the hull, but
$\pv(x_\ast)=\Ktil^{-1}\kvec(x_\ast)$ can have large entries through $\Ktil^{-1}$,
and the two pieces fail to cancel.
\end{remark}

Left unchecked, a value like $\Sigma_{x,x_\ast}\approx-4$ feeds the conditional
moments
$\pmean_{\mathrm{cond}}(x_\ast)=\pmean(x_\ast)+(\Sigma_{x,x_\ast}/\Sigma_{x,x})(\lat_{\mathrm{obs}}-\pmean(x))$
and $\pvar_{\mathrm{cond}}(x_\ast)=\Sigma_{x_\ast,x_\ast}-\Sigma_{x,x_\ast}^2/\Sigma_{x,x}$
of the distribution-aware utility (Part~III) with a regression coefficient
$\approx-100$, giving nonsensical conditional means and spurious utility peaks at
the edges of the inducing region. The fix is to never assemble the cross-block by
hand: form the joint covariance with GPyTorch and slice it,
\begin{lstlisting}
full_covar = model(torch.cat([x_samples, x_star])).covariance_matrix
Sigma      = full_covar[:N, N:]        # correct cross-block, Nystrom-consistent
\end{lstlisting}

\caveat{The shipped predict path already reads GPyTorch's joint
\texttt{covariance\_matrix} (which internally includes the conditional
correction term), so the pathology of \eqref{eq:crosscov} is confined to
\emph{hand-assembled} cross-covariances, not the production path. Any new code
that forms a cross-covariance from the manual formula outside the inducing hull
will nonetheless reintroduce it --- use the slice above.}

\subsection{Whitening}\label{sec:whitening}

The alternative reparameterisation (used by \texttt{default\_gpy},
\S\ref{sec:defaultgpy}) is \emph{Cholesky whitening}, which makes the prior a
standard normal. With the Cholesky factor $\Ktil=\Lchol\Lchol\transpose$, define
the whitened inducing values
\begin{equation}
  \widetilde{\latt}=\Lchol^{-1}(\latt-\mu_z)\;\sim\;\Gauss(0,I)\quad\text{under the prior},
\end{equation}
where $\mu_z$ is the prior mean at the inducing points. Optimisation conditions
much better in these coordinates. GPyTorch then stores the \emph{whitened}
variational parameters $(\widetilde{\vm},\widetilde{\vV})$, related to the actual
ones by
\begin{equation}
  \vm-\mu_z=\Lchol\,\widetilde{\vm},\quad \vV=\Lchol\,\widetilde{\vV}\,\Lchol\transpose;
  \qquad
  \widetilde{\vm}=\Lchol^{-1}(\vm-\mu_z),\quad \widetilde{\vV}=\Lchol^{-1}\vV\,\Lchol^{-\mathsf{T}}.
  \label{eq:whiten}
\end{equation}
The whitened mean thus encodes the \emph{deviation} from the prior mean function,
not absolute function values. Reading these stored params inside a manual moment
formula therefore requires three corrections: unwhiten the covariance
($\vV=\Lchol\widetilde{\vV}\Lchol\transpose$), unwhiten the mean
($\vm-\mu_z=\Lchol\widetilde{\vm}$), and add the mean function back,
$\pmean_\ast=\pmean(x_\ast)+\pv(x_\ast)\transpose\Lchol\widetilde{\vm}$ and
$\pvar_\ast=s(x_\ast)+\pv(x_\ast)\transpose\Lchol\widetilde{\vV}\Lchol\transpose\pv(x_\ast)$.

\begin{remark}[Zero mean in production]
Both production paths use \texttt{ZeroMean} (\texttt{gpy\_model.py}; the
eigenspace path carries no mean module), so $\mu_z=0$ and the ``add the mean
function'' step is trivially zero. The whitening maps \eqref{eq:whiten}
themselves still apply whenever the whitened \texttt{VariationalStrategy} is
selected; the mean-function machinery is retained only to guard a future non-zero
mean.
\end{remark}


\section{The eigenspace representation and the two implementations}\label{sec:eigenspace}

The two current code paths realise the \emph{same} variational GP of
\S\ref{sec:varinf} --- same prior $\Gauss(0,\Ktil)$, same arc-cosine RF kernel,
same $\Poi(\gain,\latz)$ likelihood, same expected rate \eqref{eq:fbar}, same
ELBO \eqref{eq:elbo}. They differ only in \emph{how} $q(\latt)$ is represented
and optimised:
\begin{itemize}[nosep]
  \item \texttt{vargp\_direct} (\texttt{DirectVGPModel}) projects $q$ into the
    eigenbasis of $\Ktil$ and runs a custom EM loop, using GPyTorch as a kernel
    calculator only;
  \item \texttt{default\_gpy} (\texttt{VariationalGPModel}) keeps $q$ in the full
    $\Mind$-dimensional inducing space via a stock GPyTorch
    \texttt{CholeskyVariationalDistribution} and optimises one joint ELBO.
\end{itemize}

\subsection{Eigendecomposition of the inducing gram}\label{sec:eigdecomp}

The inducing gram $\Ktil$ ($\Mind\times\Mind$) has effective rank far below
$\Mind$ (typically $\sim10$--$50$). \texttt{vargp\_direct} eigendecomposes it and
keeps only the well-conditioned directions
(\texttt{eigenspace\_utils.py:\allowbreak eigendecompose\_K\_tilde}):
\begin{equation}
\begin{gathered}
  \Ktil=\eb\,\evals\,\eb\transpose,\qquad
  \evals=\diago(\text{eigvals}_b), \\
  \eb\in\Reals^{\Mind\times\nb},\qquad \eb\transpose\eb=I_{\nb},
\end{gathered}
  \label{eq:eigdecomp}
\end{equation}
where the kept set is chosen by a \emph{relative} threshold with an absolute
floor,
\begin{equation}
  \text{keep eigenpair }j \iff \evals_j > \max\!\big(\evals_{\max}\cdot\texttt{tol},\ \texttt{tol}\big),
  \qquad \texttt{tol}=\texttt{EIGVAL\_TOL}=10^{-4}.
  \label{eq:eigthresh}
\end{equation}
\texttt{torch.linalg.eigh} returns eigenvalues in ascending order; $\eb$ collects
the kept eigenvectors as columns and $\evals$ the kept eigenvalues. The kept
count $\nb$ is \emph{dynamic}: it shifts across training iterations as the kernel
hyperparameters move, and grows by $\sim1$ when a point is appended
(\S\ref{sec:rank1}). Typical $\nb\approx10$--$11$ for the closed-loop cells, up
to $\sim50$ elsewhere. Reprojection (\S\ref{sec:recompute}) absorbs the dimension
change.

Three payoffs motivate the projection: (i)~dimensionality drops $\Mind\to\nb$, so
per-step cost falls $O(\Mind^3)\to O(\nb^3)$; (ii)~$\Ktil$ becomes
\emph{diagonal} in the new basis, $\Ktilb=\evals$, so $\Ktil^{-1}$ is the
element-wise reciprocal $\diago(1/\evals)$ --- no matrix solve; (iii)~dropping the
small eigenvalues gives implicit regularisation.

\subsection{The eigenspace state}\label{sec:eigstate}

The \texttt{DirectVariationalState} dataclass (\texttt{eigenspace\_model.py}) is
the eigenspace posterior together with all precomputed kernel projections:

\begin{center}\small
\begin{tabular}{@{}l l >{\raggedright\arraybackslash}p{8.6cm}@{}}
\toprule
field & shape & meaning \\
\midrule
\texttt{m\_b} $=\mb$ & $(\nb)$ & variational mean in eigenspace, $\mb=\eb\transpose\vm$ \\
\texttt{V\_b} $=\Vb$ & $(\nb,\nb)$ & variational covariance in eigenspace, $\Vb=\eb\transpose\vV\eb$ --- \textbf{dense} \\
\texttt{B} $=\eb$ & $(\Mind,\nb)$ & eigenvector matrix of $\Ktil$ (the basis) \\
\texttt{eigvals\_b} $=\evals$ & $(\nb)$ & kept eigenvalues of $\Ktil$ (ascending) \\
\texttt{K\_tilde\_b} $=\Ktilb$ & $(\nb,\nb)$ & inducing gram in eigenspace $=\diago(\evals)$ --- \textbf{diagonal} \\
\texttt{K\_b} & $(\Ntr,\nb)$ & cross-kernel in eigenspace $=\Kf(X,\Ipt)\,\eb$ \\
\texttt{KKtilde\_inv\_b} $=\amat$ & $(\Ntr,\nb)$ & projection $\Kf\Ktil^{-1}$ in eigenspace $=\texttt{K\_b}/\evals$ (element-wise) \\
\texttt{Kvec} $=\Kvec$ & $(\Ntr)$ & diagonal self-kernel $\Kf(x_i,x_i)$, the per-point prior variance \\
\texttt{K\_tilde} $=\Ktil$ & $(\Mind,\Mind)$ & the full inducing gram, kept only to enable the $O(\Mind)$ column-append (\S\ref{sec:rank1}) \\
\texttt{mask} & $(n_{\text{pix}})$ & RF-support pixel mask from the kernel (if \texttt{use\_mask}) \\
\bottomrule
\end{tabular}
\end{center}

\begin{remark}[Only $\Ktilb$ is diagonal]
$\Vb$ is a \emph{dense} $\nb\times\nb$ matrix even though $\Ktilb$ is diagonal;
never assume a diagonal $\Vb$. The projection row $\amat_i=(\texttt{K\_b}/\evals)_i$
is exactly the eigenspace image of $\kvec(x_i)\transpose\Ktil^{-1}$, computed by
element-wise division rather than a solve.
\end{remark}

The state is built in one pass by
\texttt{\_compute\_eigenspace\_quantities} (under \texttt{torch.no\_grad}), which
is shared by initialisation and the post-M-step recompute. Initialisation sets
the variational params to the prior:
\begin{equation}
  \mb=0,\qquad \Vb=\Ktilb \quad(\text{i.e.\ }\vm=0,\ \vV=\Ktil,\ \text{so } \KLdiv=0\text{ at start}).
\end{equation}

\subsection{Posterior moments in the eigenbasis}\label{sec:eig-moments}

With $\amat=\texttt{KKtilde\_inv\_b}=\texttt{K\_b}/\evals$, the projected moments
\eqref{eq:predmoments} at the training points become
(\texttt{eigenspace\_model.py:\allowbreak \_lambda\_moments\_eigenspace}):
\begin{align}
  \pmean &= \amat\,\mb, \notag\\
  \pvar  &= \Kvec + \diago\!\big(\amat\,(\Vb-\Ktilb)\,\amat\transpose\big)
          = \Kvec + \operatorname{rowsum}\!\big(\amat\odot(\amat\,(\Vb-\Ktilb))\big), \notag\\
  \pvar  &\leftarrow \max\!\big(\pvar,\ \texttt{LAMBDA\_VAR\_CLAMP}=10^{-6}\big),
  \label{eq:eigmoments}
\end{align}
where $\pmean,\pvar$ are the per-training-point vectors $\lambda_m,\lambda_{\mathrm{var}}$
and $\odot$ is the element-wise product; the two variance expressions coincide
because $\operatorname{rowsum}(\amat\odot(\amat M))=\diago(\amat M\amat\transpose)$
for any $M$. Term by term this equals \eqref{eq:predmoments}: $\Kvec_i=\Kf(x_i,x_i)$
(the self-kernel identity \eqref{eq:selfkernel}) and
$\amat_i(\Vb-\Ktilb)\amat_i\transpose=\kvec(x_i)\transpose\Ktil^{-1}(\vV-\Ktil)\Ktil^{-1}\kvec(x_i)$.
For \textbf{test} points the same two lines run on a freshly projected
$\amat=\Kf(x_\ast,\Ipt)\eb/\evals$ with $\Kvec=\Kf(x_\ast,x_\ast)$, so no kernel
is recomputed for training points (\texttt{predict\_eigenspace} returns
$\{\texttt{f\_pred},\texttt{lambda\_m},\texttt{lambda\_var}\}$ with
$\texttt{f\_pred}=\exp(\gain\,\lambda_m+\tfrac12\gain^2\lambda_{\mathrm{var}}+\latz)$).

\begin{remark}[Real method names]
The expected firing rate lives on
\texttt{PoissonLikelihood.\allowbreak expected\_firing\_rate(posterior)}; the
\texttt{Eigenspace\-Posterior} object exposes only \texttt{.mean} and
\texttt{.variance} (computed once at construction). Older reference docs that
place an \texttt{expected\_firing\_rate()} method on the posterior, or that name
files \texttt{estep.py}/\texttt{mstep.py}, are stale --- the live files are
\texttt{eigenspace\_estep.py}, \texttt{eigenspace\_fstep.py},
\texttt{eigenspace\_mstep.py}, \texttt{eigenspace\_training.py}.
\end{remark}

\subsection{ELBO and KL in the eigenbasis}\label{sec:eig-kl}

The eigenspace ELBO (\texttt{compute\_elbo\_eigenspace}) uses the log-likelihood
\eqref{eq:ell} and, because $\Ktilb=\diago(\evals)$, the KL \eqref{eq:kl} collapses
to sums over the kept dimensions:
\begin{align}
  \KLdiv(q\,\|\,p)
    &= \tfrac12\Big[\Tr(\Ktil^{-1}\vV)+\vm\transpose\Ktil^{-1}\vm-\nb
        +\log\lvert\Ktil\rvert-\log\lvert\vV\rvert\Big] \notag\\
    &= \tfrac12\Big[\ \textstyle\sum_j \dfrac{\Vb[j,j]}{\evals_j}
        \;+\; \sum_j \dfrac{\mb[j]^2}{\evals_j}
        \;-\; \nb
        \;+\; \sum_j \log\evals_j
        \;-\; \log\det\Vb\ \Big].
  \label{eq:eigkl}
\end{align}

\begin{remark}[Diagonal-trace assumption]
The trace term uses only $\diago(\Vb)$, which is valid \emph{only because}
$\Ktilb$ is diagonal when the ELBO is evaluated --- i.e.\ on a consistent
(kernel, basis, params) triple right after initialisation, recompute, or the
E/F-step, and before the M-step. \texttt{log\_det}$\,\Vb$ uses the full
\texttt{slogdet}; a non-positive-definite $\Vb$ warns and sets $\log\lvert\vV\rvert\!\to\!0$.
The $-\nb$ is a zero-gradient constant (cf.\ the $-\Mind$ remark after
\eqref{eq:kl}): it shifts only the reported loss magnitude.
\end{remark}

\subsection{The recompute seam after the M-step}\label{sec:recompute}

An M-step changes the kernel hyperparameters, so $\Ktil$ --- and therefore its
eigenbasis $\eb$ and eigenvalues $\evals$ --- go stale.
\texttt{DirectVGPModel.recompute\_eigenspace()} restores consistency in two steps:
\begin{enumerate}[nosep]
  \item recompute all eigenspace quantities from the \emph{new} kernel
    (\texttt{\_compute\_eigenspace\_quantities});
  \item \emph{reproject} the variational params from the old basis to the new one
    by round-tripping through full $\Mind$-space
    (\texttt{eigenspace\_utils.py:\allowbreak reproject\_variational\_params}):
    \begin{equation}
      \mb^{\text{new}}=\eb_{\text{new}}\transpose\,\eb_{\text{old}}\,\mb^{\text{old}},
      \qquad
      \Vb^{\text{new}}=\eb_{\text{new}}\transpose\big(\eb_{\text{old}}\,\Vb^{\text{old}}\,\eb_{\text{old}}\transpose\big)\eb_{\text{new}}.
      \label{eq:reproject}
    \end{equation}
\end{enumerate}
In \texttt{train\_eigenspace} this runs at the \emph{start} of every iteration
$>1$ (guarded by \texttt{n\_mstep>0 and iteration>1}), so metrics and the ELBO
\eqref{eq:eigkl} are always evaluated on a consistent triple. If $\nb$ grows, the
new dimensions carry no prior information and $\Vb^{\text{new}}$ may have tiny
eigenvalues; the subsequent E-step restores positive-definiteness.

\caveat{$\Ktilb$ is diagonal \emph{only} right after init/recompute and during
the E/F-step. Inside the M-step closure, after a hyperparameter step but before
the next recompute, the kernel has moved while the basis $\eb$ is still the
\emph{old} one, so $\Ktilb=\eb\transpose\Ktil_{\text{new}}\eb$ is \emph{not}
diagonal --- code there must use \texttt{torch.linalg.solve}, not
$\diago(1/\evals)$. Treating it as diagonal in that window is the diagonal-KL
trap (a bug fixed Jan~2025).}

\subsection{The \texttt{default\_gpy} path}\label{sec:defaultgpy}

\texttt{VariationalGPModel} (\texttt{gpy\_model.py}) realises the same model
through stock GPyTorch, keeping $q(\latt)$ in full $\Mind$-space and never
touching the eigenbasis:
\begin{itemize}[nosep]
  \item \texttt{CholeskyVariationalDistribution($\Mind$)} stores
    $q(\latt)=\Gauss(\vm,\Lchol\Lchol\transpose)$ via the Cholesky factor
    (numerically stable);
  \item \texttt{VariationalStrategy} (whitened, the default) or
    \texttt{UnwhitenedVariationalStrategy} (natural params), both with
    \texttt{jitter\_val=JITTER}$=10^{-4}$;
  \item \texttt{mean\_module=ZeroMean()}, \texttt{covar\_module=} the arc-cosine
    kernel; \texttt{forward(x)} returns the GP \emph{prior}
    $\Gauss(\text{ZeroMean}(x),\Kf(x))$, and the strategy turns that plus
    $q(\latt)$ into $q(\lat)$ at $x$.
\end{itemize}
It \emph{is} an \texttt{nn.Module}, so kernel, likelihood, and variational params
all live in one \texttt{state\_dict}. Training (\texttt{gpy\_training.py:\allowbreak train\_gpy\_default})
is a single joint ELBO maximisation over all parameters --- no EM decomposition:
\begin{equation}
\begin{aligned}
  \text{loss}=-\,&\texttt{likelihood.expected\_log\_prob}(\spk,\text{model}(X)) \\
    &+\;\texttt{model.variational\_strategy.kl\_divergence()},
\end{aligned}
\end{equation}
minimised by LBFGS (\texttt{strong\_wolfe}, \texttt{GPY\_LBFGS\_MAX\_ITER}$=20$) or
Adam. The LBFGS closure rejects a step (returns $\infty$) on out-of-bounds
kernel/likelihood params, a NaN gram, or firing-rate divergence
($\fbar$ mean$>100$ or max$>500$); after each accepted step the kernel and
likelihood are clamped (projected gradient descent). Cholesky stability comes from
\texttt{cholesky\_jitter($10^{-4}$)} with \texttt{cholesky\_max\_tries=3} (retries
at $10^{-4},10^{-3},10^{-2}$). Prediction (\texttt{gpy\_training.py:\allowbreak predict})
computes $\pmean=\texttt{posterior.mean}$, $\pvar=\max(\texttt{posterior.variance},10^{-6})$,
and $\texttt{f\_pred}=\exp(\gain\pmean+\tfrac12\gain^2\pvar+\latz)$ --- the same
\eqref{eq:fbar}, but returned under the key \texttt{lambda\_mean} rather than
\texttt{predict\_eigenspace}'s \texttt{lambda\_m}.

\subsection{Why \texttt{vargp\_direct} manages parameters directly}\label{sec:whymanage}

The whitened \texttt{VariationalStrategy} stores params relative to $\Lchol$ (the
Cholesky of $\Ktil$), making the prior $\Gauss(0,I)$; it \emph{assumes joint
autograd} optimisation, in which autograd differentiates through $\Lchol$ and
keeps the whitened coupling consistent. A closed-form Newton E-step combined with
a kernel-changing M-step (EM style) breaks that assumption. Write $L_{\text{old}}$
and $L_{\text{new}}$ for the Cholesky factor $\Lchol$ before and after such a step.
Params stored as $\vm_{\text{stored}}=L_{\text{old}}^{-1}\vm_{\text{natural}}$
become stale the moment the M-step moves $L_{\text{old}}\!\to\!L_{\text{new}}$, so
the next E-step reads $L_{\text{new}}\vm_{\text{stored}}\neq\vm_{\text{natural}}$
(empirically $\sim8\times$ errors in $\pmean$ --- the ``$L$ mismatch'',
DECISION\_LOG Q26/Q27). There are two escapes:
\begin{enumerate}[nosep]
  \item \texttt{UnwhitenedVariationalStrategy} stores natural params
    ($\Lchol$-independent), so EM is safe --- but it is slower and empirically
    less accurate (KL-gradient instability, Q29);
  \item \texttt{vargp\_direct}'s own eigenbasis whitens through $\eb/\evals$ but
    \emph{recomputes and reprojects the basis explicitly} after every M-step
    (\S\ref{sec:recompute}) instead of relying on autograd.
\end{enumerate}
This is exactly why \texttt{vargp\_direct} owns its basis and runs a custom EM
loop, while \texttt{default\_gpy} stays in the whitened single-joint-ELBO regime
that whitening was designed for. The eigenbasis $\eb$ and GPyTorch's $\Lchol$ play
analogous roles (both diagonalise/whiten the prior); the difference is who keeps
them in sync --- \texttt{vargp\_direct} does so explicitly, GPyTorch does so
implicitly via autograd (which EM defeats).

\subsection{Shared math, divergent representation}\label{sec:contrast}

\begin{table}[htbp]\centering\footnotesize
\begin{tabular}{@{}>{\raggedright\arraybackslash}p{3.0cm} >{\raggedright\arraybackslash}p{5.6cm} >{\raggedright\arraybackslash}p{5.6cm}@{}}
\toprule
aspect & \texttt{vargp\_direct} (\texttt{DirectVGPModel}) & \texttt{default\_gpy} (\texttt{VariationalGPModel}) \\
\midrule
posterior rep.\ of $q(\latt)$ & eigenspace: $\mb$ $(\nb)$, dense $\Vb$ $(\nb\!\times\!\nb)$ in basis $\eb$ of $\Ktil$ & full $\Mind$-space: \texttt{CholeskyVariational\allowbreak Distribution}, $q=\Gauss(\vm,\Lchol\Lchol\transpose)$ \\
$\Ktil^{-1}$ & trivial $\diago(1/\evals)$ & GPyTorch Cholesky of $\Ktil$ + jitter \\
whitening & own eigenbasis, \textbf{explicit reproject} after each M-step & whitened $\Gauss(0,I)$ (default) or unwhitened \\
optimiser & custom \textbf{EM}: E (Newton), F ($\gain,\latz$), M (LBFGS) $+$ recompute & single \textbf{joint ELBO} over all params (LBFGS/Adam) \\
GPyTorch role & kernel calculator only (\texttt{kernel(...).to\_dense()}) & full \texttt{ApproximateGP} $+$ \texttt{VariationalStrategy} \\
\texttt{nn.Module}? & no (plain class; dataclass state) & yes \\
KL & manual eigenspace \eqref{eq:eigkl} & \texttt{variational\_strategy.\allowbreak kl\_divergence()} \\
dim.\ reduction & $\Mind\to\nb$ eigenspace & none (full $\Mind$) \\
online rank-1 update & \texttt{extend\_model\_with\allowbreak\_new\_point} (\S\ref{sec:rank1}) & none provided \\
checkpoint & \texttt{eigenspace\_checkpoint.py} ($\mb,\Vb$; no $\eb$) & \texttt{checkpoint.py} (\texttt{state\_dict}) \\
predict return key & \texttt{f\_pred, lambda\_m, lambda\_var} & \texttt{f\_pred, lambda\_mean, lambda\_var} \\
\bottomrule
\end{tabular}
\caption{The two implementations of the identical model
(\S\ref{sec:varinf}): the ELBO, moments, likelihood, and firing rate are shared;
only the representation of $q(\latt)$ and the optimiser differ.}
\label{tab:impl-contrast}
\end{table}

Table~\ref{tab:impl-contrast} contrasts the two. Everything above the
representation line is shared: the GP prior, the arc-cosine RF kernel, the
$\Poi(\gain,\latz)$ likelihood, $\fbar$ of \eqref{eq:fbar}, the ELBO
decomposition \eqref{eq:elbo}, the divergence guards ($\fbar$ mean/max $100/500$),
the ELBO early-stopping logic, $\texttt{JITTER}=10^{-4}$, and the
inducing-points-are-training-points invariant (Remark~\ref{rem:MeqN}).

\subsection{Online growth and the checkpoint seam}\label{sec:rank1}

\paragraph{Rank-one update.}
Phase-2 active learning grows the model by one image per step
(\texttt{rank1\_update.py:\allowbreak extend\_model\_with\_new\_point}), producing a
\emph{new} \texttt{DirectVGPModel} with $\Mind{+}1$ points warm-started from the
$\Mind$-point model:
\begin{enumerate}[nosep]
  \item \textbf{Column-append to $\Ktil$.} Compute only the new column
    $\Kf(\Ipt_{\text{new}},x_{\text{new}})$ ($O(\Mind)$ kernel evals) and stitch
    it as a new row/column onto the stored \texttt{state.K\_tilde} --- this is
    \emph{why} the full $\Ktil$ is retained (\S\ref{sec:eigstate}) --- avoiding
    the $O(\Mind^2)$ full recompute. (The $(\Mind{+}1)$-eigendecomposition is
    $O(\Mind^3)$ either way.)
  \item \textbf{Eigendecompose} the new $\Ktil$ $\to$ $\eb_{\text{new}},\evals$
    ($\nb^{\text{new}}$).
  \item \textbf{Derived quantities} using the active-loop invariant $X=\Ipt$
    (so $\Kf=\Ktil$): $\texttt{K\_b}=\Ktil_{\text{new}}\eb_{\text{new}}$,
    $\Ktilb=\diago(\evals)$, $\amat=\texttt{K\_b}/\evals$; $\Kvec$ extended by the
    single new diagonal $\Kf(x_{\text{new}},x_{\text{new}})$.
  \item \textbf{Warm-start} the variational params (the legacy three-move):
    expand to full space
    $\vm_{\text{full}}=\eb_{\text{old}}\mb^{\text{old}}$,
    $\vV_{\text{full}}=\eb_{\text{old}}\Vb^{\text{old}}\eb_{\text{old}}\transpose$
    (symmetrised); pad to $\Mind{+}1$ with a new mean entry
    $\operatorname{mean}(\vm_{\text{full}})$ and a covariance
    $\vV_{\text{ext}}=I_{\Mind+1}$ with $\vV_{\text{ext}}[{:}\Mind,{:}\Mind]=\vV_{\text{full}}$
    (the new dimension gets \emph{unit variance, zero cross-covariance} --- an
    identity placeholder for the yet-unseen point); reproject into the new basis
    $\mb^{\text{new}}=\eb_{\text{new}}\transpose\vm_{\text{ext}}$,
    $\Vb^{\text{new}}=\eb_{\text{new}}\transpose\vV_{\text{ext}}\eb_{\text{new}}$.
  \item Hand the precomputed state to
    \texttt{DirectVGPModel(..., \_precomputed\_state=state)} (skipping the
    constructor's own eigenspace build). The kernel and likelihood must be
    deep-copied by the caller, since \texttt{train\_eigenspace} mutates the kernel
    in place.
\end{enumerate}
Net: one active image appends one column to $\Ktil$, bumps $\nb$ by $\sim1$, and
re-warm-starts $(\mb,\Vb)$ --- the online counterpart of the batch
initialisation.

\paragraph{Checkpoint seam.}
\texttt{eigenspace\_checkpoint.py} persists a \texttt{DirectVGPModel} (a plain
class, hence a bespoke checkpoint). It \emph{saves}
\texttt{kernel.state\_dict()}, \texttt{likelihood.state\_dict()}, the variational
params $\mb,\Vb$, \texttt{pool\_indices} (indices into the shared image pool),
pool integrity tags, and the RF-centre box. It does \emph{not} save $X$
(reconstructed as \texttt{X\_pool[pool\_indices]}, keeping checkpoints $\sim$1~KB)
nor the entire derived eigenspace $(\eb,\evals,\Ktilb,\texttt{K\_b},\amat,\Kvec)$
--- on load, \texttt{DirectVGPModel.\_\_init\_\_} recomputes the eigenspace fresh
from the loaded kernel, and the saved $\mb,\Vb$ are re-injected via
\texttt{update\_variational\_params}.

\caveat{The saved $\mb,\Vb$ are expressed in the \emph{save-time} eigenbasis
$\eb$, but are injected into a \emph{freshly recomputed} $\eb$ with \textbf{no}
reprojection --- the code assumes the recomputed eigendecomposition equals the
save-time one (``deterministic on the same hardware''). This is the documented
basis-drift fragility: under a different torch version the load-time
eigendecomposition of $\Ktil$ differs subtly (sign flips, near-threshold
eigenvectors dropping in or out of the kept set, near-degenerate subspace
rotation), so the reloaded posterior can drift by $\sim5\%$. Note the asymmetry
with the \emph{in-memory} \texttt{restore\_best} path, which \emph{does} save
$\eb$ and reprojects \eqref{eq:reproject}; only the \emph{on-disk} checkpoint
omits the basis. Integrity is guarded for the data pool (shape exact, dtype warn,
sum within tolerance), not for the basis.}

\begin{remark}[E/F/M objects only]
The custom EM loop mutates this eigenspace state --- the E-step Newton update of
$(\mb,\Vb)$ (which \texttt{detach}es both, to stop the autograd graph chaining
across iterations through $\gain$ and blowing up memory, and symmetrises $\Vb$),
the F-step on $(\gain,\latz)$, and the M-step LBFGS on the kernel
hyperparameters. The update \emph{equations} belong to the training-algorithm
section; this section fixes only the objects each step reads and writes.
\end{remark}

\endgroup  % end scoped \sloppy (see \begingroup at the top of this fragment)
